% 数学建模国赛论文模板
% Mathematical Contest in Modeling LaTeX Template

\documentclass[12pt,a4paper]{article}

% 中文支持
\usepackage{ctex}
\usepackage{fontspec}

% 数学公式
\usepackage{amsmath}
\usepackage{amssymb}
\usepackage{amsfonts}

% 图表
\usepackage{graphicx}
\usepackage{float}
\usepackage{booktabs}
\usepackage{multirow}

% 页面布局
\usepackage{geometry}
\geometry{left=2.5cm,right=2.5cm,top=2.5cm,bottom=2.5cm}

% 超链接
\usepackage{hyperref}
\hypersetup{
    colorlinks=true,
    linkcolor=blue,
    filecolor=blue,
    urlcolor=blue,
    citecolor=cyan,
}

% 代码块
\usepackage{listings}
\usepackage{xcolor}
\lstset{
    basicstyle=\ttfamily\small,
    keywordstyle=\color{blue},
    commentstyle=\color{green!60!black},
    stringstyle=\color{red},
    numbers=left,
    numberstyle=\tiny\color{gray},
    frame=single,
    breaklines=true,
    captionpos=b
}

% 算法伪代码
\usepackage{algorithm}
\usepackage{algorithmic}

% 标题信息
\title{\textbf{论文题目}}
\author{队伍名称 \\ 队员：队员1，队员2，队员3 \\ 指导教师：教师姓名}
\date{\today}

\begin{document}

% 摘要页
\maketitle
\thispagestyle{empty}

\begin{abstract}
在这里输入摘要内容。摘要应该简明扼要地说明论文的主要问题、方法、结果和结论。

\textbf{关键词：} 关键词1，关键词2，关键词3
\end{abstract}

\newpage

% 目录
\tableofcontents
\newpage

% 正文开始
\section{问题重述}

\subsection{问题背景}
在这里描述问题的背景信息。

\subsection{需要解决的问题}
在这里列出需要解决的具体问题。

\section{问题分析}

对问题进行详细分析，说明解题思路。

\section{模型假设}

列出模型的基本假设：

\begin{itemize}
    \item 假设1：描述第一个假设
    \item 假设2：描述第二个假设
    \item 假设3：描述第三个假设
\end{itemize}

\section{符号说明}

\begin{table}[H]
\centering
\caption{符号说明}
\begin{tabular}{cl}
\toprule
符号 & 说明 \\
\midrule
$x$ & 变量说明 \\
$y$ & 变量说明 \\
$\alpha$ & 参数说明 \\
\bottomrule
\end{tabular}
\end{table}

\section{模型建立与求解}

\subsection{模型一：XXX模型}

\subsubsection{模型建立}
在这里建立数学模型。

行内公式示例：$E = mc^2$

行间公式示例：
\begin{equation}
    \int_{a}^{b} f(x) dx = F(b) - F(a)
\end{equation}

\subsubsection{模型求解}
在这里求解模型，可以插入图片：

\begin{figure}[H]
    \centering
    \includegraphics[width=0.6\textwidth]{example-image.png}
    \caption{图片标题}
    \label{fig:example}
\end{figure}

\subsection{模型二：XXX模型}

继续建立和求解其他模型。

\section{模型检验}

对模型进行检验和验证。

\section{结果分析}

分析模型结果，讨论优缺点。

\section{结论与展望}

\subsection{主要结论}
总结论文的主要结论。

\subsection{模型改进与推广}
讨论模型的改进方向和推广价值。

% 参考文献
\begin{thebibliography}{99}
\bibitem{ref1} 作者1, 作者2. 文章标题[J]. 期刊名称, 年份, 卷(期): 页码.
\bibitem{ref2} 作者. 书名[M]. 出版地: 出版社, 年份.
\bibitem{ref3} 作者. 论文标题[D]. 城市: 学校名称, 年份.
\end{thebibliography}

% 附录
\newpage
\appendix
\section{附录}

\subsection{主要代码}

\begin{lstlisting}[language=Python, caption=示例代码]
import numpy as np

def example_function(x):
    """
    示例函数
    """
    return x**2 + 2*x + 1
\end{lstlisting}

\end{document}
